\newcommand{\CourseTemplatePath}{../../../../../templates/slides/}
\input{\CourseTemplatePath course_preamble.tex}
\title[Introduction to Python]{Lab: Introduction to Python}
\subtitle{Development tools, scripts, notebooks, and reproducible results}
\begin{document}
\makecoursetitle

\begin{frame}{Run a Small Program and Check Its Result}
\courseexample{Example: Convert a Temperature}\par
You want to convert a temperature from Celsius to Fahrenheit.
\begin{itemize}
\item What information does the program need as input?
\item Can you run the same calculation from a script and a notebook?
\item If the input changes from 20 to 25 degrees Celsius, what should happen to the output?
\end{itemize}
\medskip
Today's goal: run, explain, modify, and verify a small program in both file formats.
\end{frame}

\begin{frame}{What Is an IDE?}
\begin{itemize}
\item An \term{integrated development environment (IDE)} brings together tools for writing, running, and debugging programs.
\item Typical features include a file explorer, code editor, completion suggestions, debugger, terminal, and version-control tools.
\item We use \href{https://code.visualstudio.com/}{VS Code}: an extensible editor that provides an integrated Python workspace through its Python and Jupyter extensions.
\item Other options include \href{https://www.jetbrains.com/pycharm/}{PyCharm}, a Python-focused IDE, and \href{https://cursor.com/}{Cursor}, another code editor.
\end{itemize}
\medskip
The editor helps you work with code; a separately selected Python interpreter executes it.
\end{frame}

\begin{frame}{Find Your Way around VS Code}
\begin{tabularx}{\linewidth}{@{}lY@{}}
\toprule\textbf{Part of the interface} & \textbf{What you use it for}\\\midrule
Explorer & Open the project folder and find its files\\
Editor & Read and edit a script or notebook\\
Terminal & Run commands and inspect their output\\
Run and Debug & Execute code and inspect it step by step\\
Source Control & Review changes and create Git commits\\
Extensions & Add Python and Jupyter support\\\bottomrule
\end{tabularx}
\par\medskip
Open the project folder first. Then open a file and select its Python interpreter or notebook kernel.
\end{frame}

\begin{frame}{Know Which Tool Does the Work}
\small
\begin{tabularx}{\linewidth}{@{}lY@{}}
\toprule\textbf{Tool or component} & \textbf{Role}\\\midrule
Python & The language in which we express the calculation\\
Interpreter & The program that executes Python code\\
VS Code & An editor for files, projects, and notebooks\\
Package & Additional reusable code, such as NumPy or pandas\\
Environment & A project-specific Python setup and installed packages\\
Notebook kernel & The running process that executes notebook cells\\\bottomrule
\end{tabularx}
\par\medskip
Installing an editor does not install every interpreter or package. Check which Python your editor and notebook are using.
\end{frame}

\begin{frame}{Two File Formats, One Calculation}
\begin{tabularx}{\linewidth}{@{}YY@{}}
\toprule\textbf{Script: \texttt{.py}} & \textbf{Notebook: \texttt{.ipynb}}\\\midrule
Plain-text Python source & A document with code cells, text, and saved outputs\\
Run the file from the beginning & Run cells through a selected kernel\\
Useful for repeatable tasks & Useful for explanations and exploration\\
Check the printed result & Restart the kernel and run all cells to verify\\\bottomrule
\end{tabularx}
\par\medskip
Use the supplied \texttt{temperature.py} and \texttt{temperature.ipynb}. Both should produce the same converted temperature.
\end{frame}

\begin{frame}{Understand the Input, Calculation, and Output}
\courseexample{Example: Celsius to Fahrenheit}\par
\[F=\frac95 C+32.\]
\begin{itemize}
\item \textbf{Input:} a temperature \(C\) in degrees Celsius.
\item \textbf{Calculation:} multiply by \(9/5\), then add 32.
\item \textbf{Output:} the corresponding temperature \(F\) in degrees Fahrenheit.
\item \textbf{Check:} \(20\times9/5+32=68\). At \(25\) degrees Celsius, the result should be \(77\) degrees Fahrenheit.
\end{itemize}
\medskip
The starter uses Python's standard library; no additional calculation package is required.
\end{frame}

\begin{frame}{Set Up and Run the Script}
\begin{enumerate}
\item Install a supported Python 3 release, VS Code, and Microsoft's Python and Jupyter extensions.
\item Open your \texttt{computational-economics} folder in VS Code. Copy the starter files into it.
\item Use \textbf{Python: Create Environment} to create a \texttt{.venv} environment, then select its interpreter.
\item Open \texttt{temperature.py} and choose \textbf{Run Python File in Terminal}.
\item Check the displayed interpreter path, Python version, and result: 20.0 C = 68.0 F.
\end{enumerate}
\medskip
Follow the linked setup guide for your operating system. The terminal displays commands and their output; the selected interpreter executes the script.
\end{frame}

\begin{frame}{Run a Notebook from a Fresh Kernel}
\begin{enumerate}
\item Open \texttt{temperature.ipynb} and select the Python environment you created.
\item Install notebook execution support, such as \texttt{ipykernel}, in that environment when prompted.
\item Read the explanation and run the cells in order.
\item Restart the kernel, run all cells, and check that the result is still 68 degrees Fahrenheit.
\end{enumerate}
\medskip
\courseexample{Example: An Old Output Can Be Misleading}\par
A notebook can display a result saved during an earlier run. A fresh kernel clears stored variables; running all cells checks whether the saved code can recreate the result.
\end{frame}

\begin{frame}{Read the Error and Check the Setup}
\small
\begin{tabularx}{\linewidth}{@{}lY@{}}
\toprule\textbf{What you observe} & \textbf{First check}\\\midrule
Python is not found & Installation and the interpreter selected in VS Code\\
A package cannot be imported & Whether it is installed in that interpreter's environment\\
A notebook variable is undefined & Cell order; restart and run all cells\\
A file cannot be found & The project folder and the file path\\
The script runs but the answer is wrong & Inputs, calculation, units, and expected result\\\bottomrule
\end{tabularx}
\par\medskip
When asking for help, include the exact error, the relevant code, your operating system, and the selected interpreter. Check one possible cause at a time.
\end{frame}

% BEGIN MANAGED READINGS: lab_python
\begin{frame}{References and Further Reading}
\coursereading{1}
  {Prof. Feng Zhigang (2026), AI for Economic Research: Dynamic Models, Language, and Agents. [Online course]}
  {https://vonzhg.github.io/AI-ECON-2026/index.html}
  {\begin{itemize}
\item \href{https://vonzhg.github.io/AI-ECON-2026/slides/Session01_Deck5_Workbench_Python_Git_GitHub.pdf}{Session 1.5: The Workbench: Python, VS Code, Git and GitHub}
\end{itemize}}
\coursereading{2}
  {Microsoft, Visual Studio Code Documentation. [Webpages]}
  {https://code.visualstudio.com/docs}
  {\begin{itemize}
\item \href{https://code.visualstudio.com/docs/python/python-tutorial}{Getting Started with Python in VS Code}
\item \href{https://code.visualstudio.com/docs/datascience/jupyter-notebooks}{Jupyter Notebooks in VS Code}
\end{itemize}}
\coursereading{3}
  {Python Software Foundation, Python.org. [Webpage]}
  {https://www.python.org/}
  {\begin{itemize}
\item \href{https://www.python.org/downloads/}{Download Python}
\end{itemize}}
\end{frame}
% END MANAGED READINGS: lab_python
\end{document}
